% GENERATED FILE. Edit canonical lessons, then run npm run generate:books.
% canonical-source-hash: fnv1a64:da5ab8aa37e5f8c6
% canonical-lessons: PA-C220-ne-return, PA-C220-aman-ne, PA-C220-us-ne, PA-C220-ne-review, PA-C220-lia-ditta, PA-C220-ki-past, PA-C220-payoff

\chapter{Aman Bought a Book}
\label{ch:pa-a2-ne-past}
\hlchaptermodality{220}

\begin{chapteropening}
\textbf{By the end of this chapter:} \emph{I can say what Aman, or he or she, bought, gave or took, with ne after the doer.}

Say what a named person or us did to something with ne, use lia and ditta, and ask a ki question in the past.
\end{chapteropening}

\section[\texorpdfstring{\pa{ਅਮਨ} \pa{ਨੇ} \pa{ਕੰਮ} \pa{ਕੀਤਾ} … (Aman ne kamm …)}{ਅਮਨ ਨੇ ਕੰਮ ਕੀਤਾ … (Aman ne kamm …)}]{Return to the small marker}

\label{lesson:PA-C220-ne-return}



\begin{quote}
\emph{Hear:} \emph{Aman ne kamm kītā} Say \textbf{\pa{ਅਮਨ} \pa{ਨੇ} \pa{ਕੰਮ} \pa{ਕੀਤਾ}}, \textquotedblleft{}Aman did the work\textquotedblright{}.
\end{quote}

\begin{grammarlens}[title={Grammar Lens}]
After a name, the doer of a past like \textbf{\pa{ਕੀਤਾ}} takes the small marker \textbf{\pa{ਨੇ}}. After \textbf{\pa{ਮੈਂ}} it takes nothing: \textbf{\pa{ਮੈਂ} \pa{ਚਾਹ} \pa{ਪੀਤੀ}}. Hold one more verb ready: \textbf{\pa{ਖ਼ਰੀਦਣਾ}}, \textquotedblleft{}to buy\textquotedblright{}.
\end{grammarlens}

\subsection*{Your turn}
\begin{itemize}
\raggedright
  \item \emph{Hear:} \emph{Aman did the work} Say \textbf{\pa{ਅਮਨ} \pa{ਨੇ} \pa{ਕੰਮ} \pa{ਕੀਤਾ}}.
  \item \emph{Hear:} \emph{I drank tea} Say \textbf{\pa{ਮੈਂ} \pa{ਚਾਹ} \pa{ਪੀਤੀ}}.
  \item \emph{Hear:} \emph{to buy} Say \textbf{\pa{ਖ਼ਰੀਦਣਾ}}.
\end{itemize}

\begin{tcolorbox}[breakable,colback=teal!4,colframe=teal!35!black,title={Before you move on}]
Where does \textbf{\pa{ਨੇ}} go? (Right after the named doer.)
\end{tcolorbox}
\section[\texorpdfstring{\pa{ਅਮਨ} \pa{ਨੇ} \pa{ਕਿਤਾਬ} \pa{ਖ਼ਰੀਦੀ।} (Aman ne kitāb …)}{ਅਮਨ ਨੇ ਕਿਤਾਬ ਖ਼ਰੀਦੀ। (Aman ne kitāb …)}]{\pa{ਅਮਨ} \pa{ਨੇ} \pa{ਕਿਤਾਬ} \pa{ਖ਼ਰੀਦੀ} — Aman bought a book}

\label{lesson:PA-C220-aman-ne}



\begin{quote}
Say \emph{Aman did the work}. Then \emph{to buy} and \emph{a book}.
\end{quote}

\begin{grammarlens}[title={Grammar Lens: \pa{ਨੇ} beyond the work}]
The marker works with any verb that is done to something. The past still agrees with the thing: \textbf{\pa{ਕਿਤਾਬ}} is feminine, so \textbf{\pa{ਖ਼ਰੀਦੀ}}.

\begin{quote}
\textbf{\pa{ਅਮਨ} \pa{ਨੇ} \pa{ਕਿਤਾਬ} \pa{ਖ਼ਰੀਦੀ।}} — \emph{Aman ne kitāb kharīdī.} — Aman bought a book.
\end{quote}

\begin{quote}
\textbf{\pa{ਅਮਨ} \pa{ਨੇ} \pa{ਚਾਹ} \pa{ਪੀਤੀ।}} — \emph{Aman ne chāh pītī.} — Aman drank tea.
\end{quote}

\begin{quote}
\textbf{\pa{ਅਮਨ} \pa{ਨੇ} \pa{ਸੇਬ} \pa{ਖਾਧਾ।}} — \emph{Aman ne seb khādhā.} — Aman ate an apple.
\end{quote}

Going and coming are not done to anything, so they take no \textbf{\pa{ਨੇ}}: \textbf{\pa{ਅਮਨ} \pa{ਘਰ} \pa{ਗਿਆ}}.
\end{grammarlens}

\subsection*{Your turn}
\begin{itemize}
\raggedright
  \item \emph{Hear:} \emph{Aman ate a roti} Say \textbf{\pa{ਅਮਨ} \pa{ਨੇ} \pa{ਰੋਟੀ} \pa{ਖਾਧੀ।}}
  \item \emph{Hear:} \emph{Aman went to school} Say \textbf{\pa{ਅਮਨ} \pa{ਸਕੂਲ} \pa{ਗਿਆ।}}
  \item \emph{Read it:} \textbf{\pa{ਅਮਨ} \pa{ਨੇ} \pa{ਕਿਤਾਬ} \pa{ਖ਼ਰੀਦੀ।}} Point to \textbf{\pa{ਨੇ}}, then to the thing the ending agrees with.
  \item \emph{Write it:} copy \textbf{\pa{ਅਮਨ} \pa{ਨੇ} \pa{ਕਿਤਾਬ} \pa{ਖ਼ਰੀਦੀ।}} once
\end{itemize}

\begin{tcolorbox}[breakable,colback=teal!4,colframe=teal!35!black,title={Before you move on}]
Does \textbf{\pa{ਅਮਨ} \pa{ਘਰ} \pa{ਗਿਆ}} take \textbf{\pa{ਨੇ}}? (No: going is not done to anything.)
\end{tcolorbox}
\section[\texorpdfstring{\pa{ਉਸ} \pa{ਨੇ} (us ne)}{ਉਸ ਨੇ (us ne)}]{\pa{ਉਸ} \pa{ਨੇ} — he or she, before the marker}

\label{lesson:PA-C220-us-ne}



\begin{quote}
Say \emph{Aman bought a book}. Then \emph{he, she}, \textbf{\pa{ਉਹ}}, and the familiar \textbf{\pa{ਉਸ}}.
\end{quote}

\begin{grammarlens}[title={Grammar Lens: \pa{ਉਹ} turns into \pa{ਉਸ}}]
Before \textbf{\pa{ਨੇ}}, \textbf{\pa{ਉਹ}} turns into \textbf{\pa{ਉਸ}}. The verb still agrees with the thing:

\begin{quote}
\textbf{\pa{ਉਸ} \pa{ਨੇ} \pa{ਦੁੱਧ} \pa{ਪੀਤਾ।}} — \emph{us ne duddh pītā.} — He drank milk; she drank milk.
\end{quote}

\begin{quote}
\textbf{\pa{ਉਸ} \pa{ਨੇ} \pa{ਕਿਤਾਬ} \pa{ਖ਼ਰੀਦੀ।}} — \emph{us ne kitāb kharīdī.} — He bought a book; she bought a book.
\end{quote}

So \textbf{\pa{ਉਸ} \pa{ਨੇ}} can mean he or she: the sentence does not say which.
\end{grammarlens}

\subsection*{Your turn}
\begin{itemize}
\raggedright
  \item \emph{Hear:} \emph{she ate a roti} Say \textbf{\pa{ਉਸ} \pa{ਨੇ} \pa{ਰੋਟੀ} \pa{ਖਾਧੀ।}}
  \item \emph{Hear:} \emph{he drank tea} Say \textbf{\pa{ਉਸ} \pa{ਨੇ} \pa{ਚਾਹ} \pa{ਪੀਤੀ।}}
  \item \emph{Read it:} \textbf{\pa{ਉਸ} \pa{ਨੇ} \pa{ਦੁੱਧ} \pa{ਪੀਤਾ।}} Give its meaning.
  \item \emph{Write it:} copy \textbf{\pa{ਉਸ} \pa{ਨੇ}} once
\end{itemize}

\begin{tcolorbox}[breakable,colback=teal!4,colframe=teal!35!black,title={Before you move on}]
\emph{He} before \textbf{\pa{ਨੇ}}? (\textbf{\pa{ਉਸ}}.) \emph{He went}? (\textbf{\pa{ਉਹ} \pa{ਗਿਆ}}, with no \textbf{\pa{ਨੇ}}.)
\end{tcolorbox}
\section[\texorpdfstring{\pa{ਅਮਨ} \pa{ਨੇ} · \pa{ਉਸ} \pa{ਨੇ} (Aman ne · us ne)}{ਅਮਨ ਨੇ · ਉਸ ਨੇ (Aman ne · us ne)}]{Hear and say the doer again}

\label{lesson:PA-C220-ne-review}



\begin{quote}
\emph{Hear:} \emph{Aman ne kitāb kharīdī} Say \textbf{\pa{ਅਮਨ} \pa{ਨੇ} \pa{ਕਿਤਾਬ} \pa{ਖ਼ਰੀਦੀ।}}
\end{quote}

\subsection*{Your turn}
\begin{itemize}
\raggedright
  \item \emph{Hear:} \emph{she bought a book} Say \textbf{\pa{ਉਸ} \pa{ਨੇ} \pa{ਕਿਤਾਬ} \pa{ਖ਼ਰੀਦੀ।}}
  \item \emph{Hear:} \emph{Aman ate an apple} Say \textbf{\pa{ਅਮਨ} \pa{ਨੇ} \pa{ਸੇਬ} \pa{ਖਾਧਾ।}}
  \item \emph{Hear:} \emph{he went home} Say \textbf{\pa{ਉਹ} \pa{ਘਰ} \pa{ਗਿਆ।}}
  \item \emph{Read it:} \textbf{\pa{ਉਸ} \pa{ਨੇ} \pa{ਚਾਹ} \pa{ਪੀਤੀ।}} Give its meaning.
  \item \emph{Write it:} cover the model and write \textbf{\pa{ਉਸ} \pa{ਨੇ} \pa{ਦੁੱਧ} \pa{ਪੀਤਾ।}} once
\end{itemize}

\begin{tcolorbox}[breakable,colback=teal!4,colframe=teal!35!black,title={Before you move on}]
Which of the three takes no \textbf{\pa{ਨੇ}}? (Going home.)
\end{tcolorbox}
\section[\texorpdfstring{\pa{ਲਿਆ} · \pa{ਦਿੱਤਾ} (liā · dittā)}{ਲਿਆ · ਦਿੱਤਾ (liā · dittā)}]{\pa{ਲਿਆ}, \pa{ਦਿੱਤਾ} — took, gave}

\label{lesson:PA-C220-lia-ditta}



\begin{quote}
Say \emph{to take} and \emph{to give}: \textbf{\pa{ਲੈਣਾ}}, \textbf{\pa{ਦੇਣਾ}}. Then \emph{Aman bought a book}.
\end{quote}

\begin{grammarlens}[title={Grammar Lens: two short pasts}]
\textbf{\pa{ਲੈਣਾ}} and \textbf{\pa{ਦੇਣਾ}} have short pasts of their own: \textbf{\pa{ਲਿਆ}}, \textbf{\pa{ਦਿੱਤਾ}}. They agree with the thing, like every past you have met:

\begin{quote}
\textbf{\pa{ਮੈਂ} \pa{ਪਾਣੀ} \pa{ਲਿਆ।}} — \emph{maiṁ pāṇī liā.} — I took water.
\end{quote}

\begin{quote}
\textbf{\pa{ਅਮਨ} \pa{ਨੇ} \pa{ਪੈਸਾ} \pa{ਦਿੱਤਾ।}} — \emph{Aman ne paisā dittā.} — Aman gave money.
\end{quote}

\begin{quote}
\textbf{\pa{ਮੈਂ} \pa{ਅਮਨ} \pa{ਨੂੰ} \pa{ਕਿਤਾਬ} \pa{ਦਿੱਤੀ।}} — \emph{maiṁ Aman nūṁ kitāb dittī.} — I gave Aman a book.
\end{quote}

The one who receives takes \textbf{\pa{ਨੂੰ}}. CDIAL lists Punjabi \emph{dittā}, \textquotedblleft{}given\textquotedblright{} (CDIAL 6140).
\end{grammarlens}

\subsection*{Your turn}
\begin{itemize}
\raggedright
  \item \emph{Hear:} \emph{she gave the book} Say \textbf{\pa{ਉਸ} \pa{ਨੇ} \pa{ਕਿਤਾਬ} \pa{ਦਿੱਤੀ।}}
  \item \emph{Hear:} \emph{I took milk} Say \textbf{\pa{ਮੈਂ} \pa{ਦੁੱਧ} \pa{ਲਿਆ।}}
  \item \emph{Read it:} \textbf{\pa{ਮੈਂ} \pa{ਅਮਨ} \pa{ਨੂੰ} \pa{ਕਿਤਾਬ} \pa{ਦਿੱਤੀ।}} Give its meaning.
  \item \emph{Write it:} copy \textbf{\pa{ਲਿਆ}} and \textbf{\pa{ਦਿੱਤਾ}} once
\end{itemize}

\begin{tcolorbox}[breakable,colback=teal!4,colframe=teal!35!black,title={Before you move on}]
\emph{Took}? (\textbf{\pa{ਲਿਆ}}.) \emph{Gave}, a book? (\textbf{\pa{ਦਿੱਤੀ}}.)
\end{tcolorbox}
\section[\texorpdfstring{\pa{ਕੀ} \pa{ਤੁਸੀਂ} \pa{ਰੋਟੀ} \pa{ਖਾਧੀ}? (kī tusīṁ roṭī khādhī?)}{ਕੀ ਤੁਸੀਂ ਰੋਟੀ ਖਾਧੀ? (kī tusīṁ roṭī khādhī?)}]{\pa{ਕੀ} \pa{ਤੁਸੀਂ} \pa{ਰੋਟੀ} \pa{ਖਾਧੀ}? — did you eat the roti?}

\label{lesson:PA-C220-ki-past}



\begin{quote}
Ask \emph{do you drink tea?} Then say \emph{you ate a roti}.
\end{quote}

\begin{grammarlens}[title={Grammar Lens: \pa{ਕੀ} in the past}]
\textbf{\pa{ਕੀ}} turns a past sentence into a question the same way it turns a present one:

\begin{quote}
\textbf{\pa{ਕੀ} \pa{ਤੁਸੀਂ} \pa{ਰੋਟੀ} \pa{ਖਾਧੀ}?} — \emph{kī tusīṁ roṭī khādhī?} — Did you eat the roti?
\end{quote}

\begin{quote}
\textbf{\pa{ਹਾਂ} \pa{ਜੀ}, \pa{ਮੈਂ} \pa{ਰੋਟੀ} \pa{ਖਾਧੀ।}} — \emph{hāṁ jī, maiṁ roṭī khādhī.} — Yes, I ate the roti.
\end{quote}

\begin{quote}
\textbf{\pa{ਨਹੀਂ} \pa{ਜੀ}, \pa{ਮੈਂ} \pa{ਅਜੇ} \pa{ਰੋਟੀ} \pa{ਨਹੀਂ} \pa{ਖਾਧੀ।}} — \emph{nahīṁ jī, maiṁ aje roṭī nahīṁ khādhī.} — No, I haven't eaten the roti yet.
\end{quote}
\end{grammarlens}

\subsection*{Your turn}
\begin{itemize}
\raggedright
  \item \emph{Hear:} \emph{did you drink tea?} Say \textbf{\pa{ਕੀ} \pa{ਤੁਸੀਂ} \pa{ਚਾਹ} \pa{ਪੀਤੀ}?}
  \item \emph{Hear:} \emph{did Aman buy the book?} Say \textbf{\pa{ਕੀ} \pa{ਅਮਨ} \pa{ਨੇ} \pa{ਕਿਤਾਬ} \pa{ਖ਼ਰੀਦੀ}?}
  \item \emph{Read it:} \textbf{\pa{ਨਹੀਂ} \pa{ਜੀ}, \pa{ਮੈਂ} \pa{ਅਜੇ} \pa{ਰੋਟੀ} \pa{ਨਹੀਂ} \pa{ਖਾਧੀ।}} Give its meaning.
  \item \emph{Write it:} copy \textbf{\pa{ਕੀ} \pa{ਤੁਸੀਂ} \pa{ਰੋਟੀ} \pa{ਖਾਧੀ}?} once
\end{itemize}

\begin{tcolorbox}[breakable,colback=teal!4,colframe=teal!35!black,title={Before you move on}]
Ask a guest whether they ate the roti. Then answer as the guest.
\end{tcolorbox}
\section[\texorpdfstring{\pa{ਅਮਨ} \pa{ਨੇ} \pa{ਕਿਤਾਬ} \pa{ਖ਼ਰੀਦੀ।} (Aman ne kitāb …)}{ਅਮਨ ਨੇ ਕਿਤਾਬ ਖ਼ਰੀਦੀ। (Aman ne kitāb …)}]{Who did what: a four-skill payoff}

\label{lesson:PA-C220-payoff}



\begin{quote}
Say \emph{Aman bought a book}, then \emph{she drank tea}.
\end{quote}

\subsection*{Your turn: four skills}
\begin{itemize}
\raggedright
  \item \emph{Hear:} \emph{Aman gave money} Say \textbf{\pa{ਅਮਨ} \pa{ਨੇ} \pa{ਪੈਸਾ} \pa{ਦਿੱਤਾ।}}
  \item \emph{Hear:} \emph{did he drink milk?} Say \textbf{\pa{ਕੀ} \pa{ਉਸ} \pa{ਨੇ} \pa{ਦੁੱਧ} \pa{ਪੀਤਾ}?}
  \item \emph{Read it:} \textbf{\pa{ਉਸ} \pa{ਨੇ} \pa{ਕਿਤਾਬ} \pa{ਲਈ} \pa{ਪੈਸਾ} \pa{ਦਿੱਤਾ।}} Give its meaning: he paid money for the book.
  \item \emph{Write it:} cover the line and write \textbf{\pa{ਅਮਨ} \pa{ਨੇ} \pa{ਕਿਤਾਬ} \pa{ਖ਼ਰੀਦੀ।}}
  \item Check: \textbf{\pa{ਨੇ}} after the named doer, and the ending after the thing.
\end{itemize}

\begin{tcolorbox}[breakable,colback=teal!4,colframe=teal!35!black,title={Before you move on}]
Tell what Aman did yesterday: he went to the market, bought a book, and drank tea.
\end{tcolorbox}
